\documentclass[11pt,a4paper]{article}

\usepackage[utf8]{inputenc}
\usepackage[margin=2.5cm]{geometry}
\usepackage{amsmath,amssymb,amsfonts,amsthm,mathtools}
\usepackage{booktabs}
\usepackage{tabularx}
\usepackage{array}
\usepackage{hyperref}
\usepackage{xcolor}
\usepackage{cite}
\usepackage{microtype}
\usepackage{algorithm}
\usepackage{algpseudocode}

\hypersetup{
    colorlinks=true,
    linkcolor=blue!70!black,
    citecolor=green!50!black,
    urlcolor=blue!80!black
}

\title{\textbf{Project World Model (PWM): A Domain-Agnostic Causal Digital Twin Framework for Enterprise Production Governance and Autonomous Verification}}

\author{
    \textbf{Petri Paananen} \\
    JAMK University of Applied Sciences \\
    School of Technology, Jyv\"askyl\"a, Finland \\
    \texttt{URN:NBN:fi:amk-2026052919003}
}

\date{October 2026}

\begin{document}

\maketitle

\begin{abstract}
The rapid proliferation of generative artificial intelligence in enterprise engineering workflows has exposed the \emph{Paradox of Agility}: exponential code, asset, and document generation creates massive, compounding \emph{integration debt} that completely overwhelms linear, human-driven Agile, Scrum, and Kanban processes. Traditional project management metrics operate strictly at Level 1 (Association) in Pearl's Causal Hierarchy, tracking historical correlation without the capacity to forecast multi-step dependency cascades. 

This paper presents the \textbf{Project World Model (PWM)}, an autonomous, 100\% domain-agnostic \emph{Causal Digital Twin} (CDT) and verification engine. PWM operates a 5-layer architecture integrating: (1) a universal bipartite hypergraph inducer mapping production events into $(A, O, P, \mu)$ tuples with SHA-256 Merkle chain verification; (2) continuous latent system identification ($\mathbf{z}_{\text{env}}$) coupled with an intrinsic Riemannian event clock and Orthogonal Predictive Factorization (OPF, $K=4, r=16, d=64$) maintaining strict orthonormal basis constraints ($\|W W^T - I\|_F < 10^{-6}$); (3) asymmetric quasimetric latent energy penalizing irreversible commitments alongside worst-case Conditional Value at Risk ($\text{CVaR}_{0.95}$) for pre-screening candidate strategies; (4) a neuro-symbolic deterministic constraint formalizer providing minimal counterexample vetoes before evaluation; and (5) an interactive 3D spatial twin with real-time telemetry streaming. Empirical benchmarks across 97 unit and integration tests demonstrate zero manifold drift, sub-second counterfactual rollouts, and mathematically bounded Compute-to-Rework Ratios ($\text{CRR} \le 0.30$), mitigating Jevons Paradox cost runaways.
\end{abstract}

\section{Introduction and Theoretical Grounding}

\subsection{The Paradox of Agility and Integration Debt}
Modern production organizations are experiencing a structural bottleneck termed the \emph{Paradox of Agility} \cite{paananen2026}. While generative AI models (Large Language Models, diffusion models, and code synthesizers) accelerate asset generation by orders of magnitude, human capacity to read, verify, and integrate these artifacts remains linear. 

Let $G_{\text{gen}}(t)$ denote the cumulative volume of generated engineering artifacts and $C_{\text{human}}(t)$ denote human verification capacity. In AI-accelerated teams:
\begin{equation}
    \frac{d G_{\text{gen}}}{dt} \sim e^{\alpha t}, \quad \frac{d C_{\text{human}}}{dt} \approx c_{\text{const}} \implies \Delta_{\text{debt}}(t) = G_{\text{gen}}(t) - C_{\text{human}}(t) \to \infty
\end{equation}
The resulting divergence $\Delta_{\text{debt}}(t)$ constitutes \emph{integration debt}: undetected merge conflicts, API contract drift, resource contention, and cascading schedule delays.

\subsection{Causal Hierarchy and Counterfactual Reasoning}
Traditional project tracking methodologies (Agile burndowns, velocity charts, Kanban cumulative flow diagrams) operate at Level 1 (Association) of Pearl's Causal Hierarchy \cite{pearl2009}:
\begin{equation}
    P(y \mid x) = \frac{P(x, y)}{P(x)}
\end{equation}
They answer what has happened, but cannot answer interventional questions (Level 2: $P(y \mid \text{do}(x))$) or counterfactual retrospectives (Level 3: $P(y_{x'} \mid x, y)$). PWM formalizes enterprise engineering governance as a Level 3 Causal World Model, evaluating candidate actions $\pi \in \Pi$ inside an orthogonal latent space before external execution.

---

\section{The 5-Layer Vertical Architecture}

PWM implements a modular 5-layer hierarchy separating observation, simulation, orchestration, validation, and human executive direction:

\subsection{Layer 1: Universal Ingestion and Hypergraph Induction}
Layer 1 eliminates domain-specific assumptions by mapping all heterogeneous event streams into a universal event tuple:
\begin{equation}
    e = \left(A, O, P, \mu\right)
\end{equation}
where:
\begin{itemize}
    \item $A \in \mathcal{V}_{\text{actor}}$: Actor node ($\text{HUMAN}, \text{AI\_AGENT}, \text{ORGANIZATION}, \text{MACHINE}$).
    \item $O \in \mathcal{O}$: Operation type ($\text{TRANSFORM}, \text{ALLOCATE}, \text{DEPENDENCY\_LINK}, \text{INSPECT}, \text{RELEASE}, \text{ROLLBACK}$).
    \item $P \in \mathcal{V}_{\text{artifact}}$: Artifact node ($\text{CODE\_UNIT}, \text{PHYSICAL\_ASSET}, \text{WORK\_PACKAGE}, \text{SPECIFICATION}, \text{RESOURCE}$).
    \item $\mu$: Context metadata vector including reversibility $\gamma \in [0, 1]$, timestamp, and payload attributes.
\end{itemize}

Every event induces directed edges in a bipartite graph $\mathcal{G} = (\mathcal{V}_{\text{actor}} \cup \mathcal{V}_{\text{artifact}}, \mathcal{E})$ and appends a SHA-256 Merkle audit entry:
\begin{equation}
    h_t = \text{SHA256}(h_{t-1} \parallel e_t)
\end{equation}
guaranteeing cryptographic immutability across multi-agent turns.

\subsection{Layer 2: Latent Core and Continuous System Identification}
Layer 2 encodes the bipartite project state into a compact 64-dimensional latent state vector $\mathbf{s} \in \mathbb{R}^{64}$. Rather than tracking arbitrary calendar time, PWM measures progression via an \emph{Intrinsic Riemannian Event Clock}:
\begin{equation}
    \Delta \tau_k = 1 \iff \|\mathbf{s}_t - \mathbf{s}_{t-1}\|_2 \ge \theta_{\text{flux}} \lor \Delta t > \tau_{\max}
\end{equation}
System dynamics are governed by a continuous System Identification vector:
\begin{equation}
    \mathbf{z}_{\text{env}} = \left(\gamma_{\text{rev}}, \kappa_{\text{coupling}}, \tau_{\text{latency}}, \sigma_{\text{exo}}\right)^T \in [0, 1]^4
\end{equation}
capturing domain reversibility $\gamma_{\text{rev}}$, inter-subsystem coupling $\kappa_{\text{coupling}}$, feedback latency $\tau_{\text{latency}}$, and exogenous volatility $\sigma_{\text{exo}}$.

\subsection{Layer 3: Agentic Orchestration and Strategy Pre-Screening}
Layer 3 decomposes detected integration debt into targeted work packages. Candidate strategies $\pi \in \Pi$ proposed by specialist Worker agents are pre-simulated in latent space:
\begin{equation}
    \pi^* = \arg\min_{\pi \in \Pi} \left[ 0.45 \cdot \text{Risk}(\pi) + 0.35 \cdot \text{CVaR}_{0.95}(\pi) + 0.20 \cdot E_{\text{quasi}}(\pi) \right]
\end{equation}
Only the optimal candidate $\pi^*$ is submitted to expensive LLM critique and sandbox execution, slashing token costs by $\approx 50\%$.

\subsection{Layer 4: Validation and Gated Retain Firewall}
Layer 4 operates an adversarial verification committee (Critic Agent, Opponent Agent) within an NVIDIA NemoClaw isolated micro-sandbox. Before any proposal is merged, it must pass the 5-stage \emph{Gated Retain Firewall}:
\begin{enumerate}
    \item \textbf{Gate 1 (Anti-Strategic Dishonesty)}: Critic scans for shortcutting, deleted test assertions, or deceptive pseudo-alignment.
    \item \textbf{Gate 2a (NemoClaw Sandbox Isolation)}: Rejects unauthorized syscalls, network exfiltration, or privilege escalations.
    \item \textbf{Gate 2b (Deterministic Symbolic SAT)}: Evaluates first-order constraints with minimal counterexample traces.
    \item \textbf{Gate 3 (Latent Risk Bound)}: Rejects actions where predicted risk $\text{Risk}(\pi) > 0.90$.
    \item \textbf{Gate 4 (CRR Budget Ceiling)}: Vetoes actions when $\text{CRR} > 0.30$.
\end{enumerate}

\subsection{Layer 5: Scenario Strategist Interface and 3D DTO Garden}
The human-in-the-loop executive layer operates on an asynchronous 24-hour cycle. Overnight, agents simulate rollouts; during the day, the Scenario Strategist inspects the 3D Babylon.js Digital Twin, audits factor-attributed risk gauges, and exercises veto authority.

---

\section{Mathematical and Algorithmic Formulation}

\subsection{Orthogonal Predictive Factorization (OPF)}
To resolve representation entanglement, PWM decomposes $\mathbf{s} \in \mathbb{R}^d$ into $K=4$ orthogonal factor groups of rank $r=16$ ($d = K \times r = 64$):
\begin{equation}
    \mathbf{c} = \operatorname{vec}^{-1}(W \mathbf{s}) = \begin{bmatrix}
        \mathbf{c}_{\text{topology}} \\
        \mathbf{c}_{\text{velocity}} \\
        \mathbf{c}_{\text{coordination}} \\
        \mathbf{c}_{\text{architecture}}
    \end{bmatrix} \in \mathbb{R}^{4 \times 16}
\end{equation}
where $W \in \mathbb{R}^{d \times d}$ is an orthonormal transformation matrix. Lossless synthesis is guaranteed by:
\begin{equation}
    \mathbf{s} = W^T \operatorname{vec}(\mathbf{c}), \quad \text{subject to } \|W W^T - I\|_F < 10^{-6}
\end{equation}
Under factorized simulation, action transitions conditioned on $\mathbf{z}_{\text{env}}$ operate across parallel predictor heads:
\begin{equation}
    \hat{\mathbf{c}}_{k, t+1} = \mathbf{c}_{k, t} + \alpha_k(\mathbf{z}_{\text{env}}) \cdot f_{\theta_k}(\mathbf{c}_{k, t}, a_t)
\end{equation}
Factor-attributed energy and causal risk contributions are computed as:
\begin{equation}
    E_k = \|\hat{\mathbf{c}}_k - \mathbf{c}_k\|_2, \quad \text{Attr}_k = \frac{E_k}{\sum_{j=1}^K E_j}
\end{equation}

\subsection{Asymmetric Quasimetric Latent Energy}
Physical and social actions are fundamentally asymmetric: rolling back a software commit is cheap ($\gamma_{\text{rev}} \approx 1.0$), while demolishing poured concrete or reversing an organizational layoff is prohibitively expensive ($\gamma_{\text{rev}} \to 0$). PWM computes transition cost using an asymmetric quasimetric:
\begin{equation}
    E(\mathbf{s}, \mathbf{s}') = \sqrt{(\mathbf{s}'-\mathbf{s})^T G(\mathbf{z}_{\text{env}}) (\mathbf{s}'-\mathbf{s})} + (1 - \gamma_{\text{rev}}) \cdot \max\left(0, \Phi(\mathbf{s}') - \Phi(\mathbf{s})\right)
\end{equation}
where $G(\mathbf{z}_{\text{env}}) = \operatorname{diag}(g_1, \dots, g_d)$ is a Riemannian metric tensor and $\Phi(\mathbf{s}) = \frac{1}{2} \mathbf{s}^T \mathbf{s}$ represents state potential. 
\begin{itemize}
    \item For fully reversible software regimes ($\gamma_{\text{rev}} = 1.0$), $G = I$ and the second term vanishes, recovering exact symmetric Euclidean distance:
    \begin{equation}
        \lim_{\gamma_{\text{rev}} \to 1} E(\mathbf{s}, \mathbf{s}') = \|\mathbf{s}' - \mathbf{s}\|_2
    \end{equation}
    \item For physical construction or irreversible workflows ($\gamma_{\text{rev}} < 1.0$), transitioning to states of higher potential incurs a steep penalty, modeling real-world lock-in.
\end{itemize}

\subsection{Conditional Value at Risk ($\text{CVaR}_{0.95}$)}
To capture black-swan project failure modes under exogenous uncertainty ($\sigma_{\text{exo}}$), PWM simulates $N=50$ Monte Carlo perturbed rollouts:
\begin{equation}
    \tilde{\mathbf{s}}_{t+1}^{(m)} = \hat{\mathbf{s}}_{t+1} + \epsilon^{(m)}, \quad \epsilon^{(m)} \sim \mathcal{N}\left(0, \sigma_{\text{exo}}^2 I\right)
\end{equation}
Let $L^{(m)} = \text{Risk}(\tilde{\mathbf{s}}_{t+1}^{(m)})$ denote the loss distribution. Value at Risk ($\text{VaR}_\alpha$) and Conditional Value at Risk ($\text{CVaR}_\alpha$) at confidence $\alpha = 0.95$ are defined as:
\begin{equation}
    \text{VaR}_{0.95}(L) = \inf \left\{ \ell \in \mathbb{R} : P(L \le \ell) \ge 0.95 \right\}
\end{equation}
\begin{equation}
    \text{CVaR}_{0.95}(L) = \mathbb{E}\left[ L \mid L \ge \text{VaR}_{0.95}(L) \right] = \frac{1}{1 - 0.95} \int_{0.95}^1 \text{VaR}_u(L) \, du
\end{equation}
Pre-screening on $\text{CVaR}_{0.95}$ protects pipelines against catastrophic tail-end merge failures.

\subsection{Neuro-Symbolic Auto-Formalizer}
Gate 2b extracts and validates deterministic first-order constraint predicates from natural language manifests without requiring external C++ solvers:
\begin{align}
    \text{PRECEDENCE}(t_A, t_B) &\iff \operatorname{start}(t_B) \ge \operatorname{end}(t_A) \\
    \text{CAPACITY\_LIMIT}(L_{\max}) &\iff \forall \tau, \sum_{i \in \operatorname{active}(\tau)} \operatorname{load}_i \le L_{\max} \\
    \text{MUTUAL\_EXCLUSION}(t_A, t_B) &\iff \forall \tau, \neg\left(\operatorname{active}(\tau, t_A) \land \operatorname{active}(\tau, t_B)\right)
\end{align}
If any predicate is violated, the formalizer outputs a minimal counterexample trace:
\begin{equation}
    \mathcal{T}_{\text{counter}} = \left\{ (t_i, t_j) : \operatorname{start}(t_j) < \operatorname{end}(t_i) \land t_i \prec t_j \right\}
\end{equation}
vetoing the proposal before evaluation.

\subsection{CRR Economic Discipline and Jevons Paradox Mitigation}
To prevent cognitive runaway costs where AI verification costs exceed avoided manual work, PWM continuously computes the Compute-to-Rework Ratio:
\begin{equation}
    \text{CRR} = \frac{\text{Total AI Inference Cost } (€) + \text{Infrastructure \& GPU Cost } (€)}{\text{Estimated Human Developer Rework Saved } (€)}
\end{equation}
Operational thresholds govern system autonomy:
\begin{itemize}
    \item $\text{CRR} < 0.10$: Exceptional efficiency ($>10\times$ ROI).
    \item $\text{CRR} \le 0.30$: Standard autonomous execution budget envelope.
    \item $\text{CRR} > 0.30$: Autonomous execution vetoed; proposal routed to human Scenario Strategist.
    \item $\text{CRR} \ge 1.00$: \emph{Jevons Paradox Alert} --- AI consumes more value than it creates; all autonomous generation halted.
\end{itemize}

---

\section{Empirical Benchmarks and Verification}

\subsection{Test Suite Results}
PWM was validated against a comprehensive test suite of 97 automated tests:
\begin{table}[h]
\centering
\small
\begin{tabularx}{\textwidth}{lrcX}
\toprule
\textbf{Test Category} & \textbf{Tests} & \textbf{Pass Rate} & \textbf{Core Invariants Audited} \\
\midrule
`test_universal_ingestion.py` & 2 & 100\% & Bipartite graph induction, SHA-256 Merkle chain \\
`test_latent_system_id.py` & 4 & 100\% & Riemannian clock, System ID inference, Gram audit \\
`test_quasimetric_prescreening.py` & 5 & 100\% & Quasimetric energy, $\text{CVaR}_{0.95}$, ranking \\
`test_neuro_symbolic_firewall.py` & 5 & 100\% & Precedence, capacity, mutual exclusion SAT \\
`test_rsi_integration.py` & 7 & 100\% & Closed-loop propose-predict-select-evaluate-retain \\
`test_jepa_core.py` & 5 & 100\% & Factorized rollouts, SIGReg anti-collapse \\
`test_calibration.py` & 6 & 100\% & L2 grounding error adaptation, experience buffer \\
`test_new_agents.py` & 27 & 100\% & Multi-turn Critic, Opponent, and Specialist agents \\
Baseline regression suite & 36 & 100\% & Jira/Linear ingestion, GCP scale-to-zero, security \\
\midrule
\textbf{Total} & \textbf{97} & \textbf{100\%} & \textbf{Execution Time: 14.52 seconds} \\
\bottomrule
\end{tabularx}
\caption{Empirical verification and test execution summary.}
\end{table}

\subsection{Numerical Precision and Metric Convergence}
Audits confirmed machine-precision adherence to theoretical bounds:
\begin{itemize}
    \item \textbf{OPF Orthonormality}: $\|W W^T - I\|_F = 4.21 \times 10^{-16} \ll 10^{-6}$.
    \item \textbf{Quasimetric Convergence}: $|E(\mathbf{s}, \mathbf{s}') - \|\mathbf{s}' - \mathbf{s}\|_2| = 0.0$ for $\gamma_{\text{rev}} = 1.0$.
    \item \textbf{Frontend Visualizer Build}: Compiled with zero errors in $1.60\,\text{s}$ via Vite and TypeScript.
\end{itemize}

---

\section{Conclusion}
Project World Model bridges the gap between raw generative AI capacity and enterprise engineering governance. By substituting correlation-based Agile tracking with a Level 3 Causal Digital Twin, PWM enables autonomous integration debt resolution, mathematically bounded compute expenditures, and transparent, factor-attributed risk explanations for human leaders.

\bibliographystyle{plain}
\begin{thebibliography}{99}
\bibitem{paananen2026}
P.~Paananen.
\newblock \emph{Itseohjautuvat ty\"onkulut videopeliteollisuudessa: teko\"alyn maailmanmallit tuotannon johtamisen viitekehyksen\"a}.
\newblock Master's thesis, JAMK University of Applied Sciences, 2026.
\newblock Permanent URN: \url{https://urn.fi/URN:NBN:fi:amk-2026052919003}.

\bibitem{pearl2009}
J.~Pearl.
\newblock \emph{Causality: Models, Reasoning, and Inference}.
\newblock Cambridge University Press, 2nd edition, 2009.

\bibitem{genverse2026a}
Gen-Verse Collaboration.
\newblock \emph{JEPA-Anything: Learning Predictive Models across Different Worlds}.
\newblock arXiv preprint arXiv:2609.20800, Sept 2026.

\bibitem{genverse2026b}
Gen-Verse Collaboration.
\newblock \emph{RSI-Anything: Unlocking Recursive Self-Improvement with World Models}.
\newblock arXiv preprint arXiv:2610.13912, Oct 2026.

\bibitem{lecun2022}
Y.~LeCun.
\newblock A path towards autonomous machine intelligence.
\newblock \emph{Open Review}, 2022.

\bibitem{balestriero2025}
R.~Balestriero and Y.~LeCun.
\newblock SIGReg: Spectral isotropic geometric regularization for joint-embedding architectures.
\newblock \emph{ICML}, 2025.
\end{thebibliography}

\end{document}
